\section{The compressed algorithm and its bit complexity}\label{sec:algorithm}

\subsection{One finite solve and a compact exact answer}

The implementation in \code{fastunknot.twist.tail} has a deliberately small
mathematical interface. It validates the run list, selects a specified run
or the largest run, and computes \([a,c]\) and \(m_0\). At or below the threshold
it uses the ordinary macro backend. Above the threshold it changes
only the selected absolute exponent to \(m_0\), computes the full reference
homology and boundary ranks, and applies \cref{thm:recurrence}.

The returned degree profile consists of a dictionary of exceptional
degree--dimension pairs and at most one interval
\[
(\ell,u,\eta),\qquad
\text{meaning }\beta_h=\eta\text{ for every integer }\ell\le h\le u.
\]
The exceptional degrees are disjoint from that interval. Endpoints are
inclusive. A single-degree query uses these two components directly.
The total rank is the sum of the point coefficients plus
\((u-\ell+1)\eta\) for each interval. An explicit expansion routine is
available, but has a caller-visible limit on the number of emitted terms.

\begin{center}
\begin{minipage}{0.94\linewidth}
\begin{lstlisting}
validate the original braid and retain its component count
choose the distinguished run, with sign epsilon and length m
compute context endpoints a, c and threshold m0 = c - a + 2
if m <= m0:
    return complete ordinary macro homology as a point profile
reference = the same signed runs with selected length m0
solve the complete reference complex
hc = c + 1 if epsilon > 0 else a - 2
M = reference.chain_dimensions[hc]
r = reference.boundary_ranks[hc]
eta = M - 2*r
shift the appropriate reference boundary region by m - m0
insert the proved constant interval, if it is nonempty
return exact rank, compact profile, original components,
       reference output and recurrence certificate
\end{lstlisting}
\end{minipage}
\end{center}

The certificate records the threshold, degree endpoints, selected position,
sign, reference dimensions, central boundary rank, slope, and finite solve.
It is a transparent computational witness. The implementation does not
claim that arbitrary certificate JSON is a formal proof of the topological
normal-form theorem. The independent audit reconstructs the recurrence
without calling the production recurrence routine.

\subsection{A precise parameterized guarantee}

\begin{theorem}[Cost in the context and exponent]\label{thm:bit-cost}
Let a fixed-context family be as in \cref{thm:recurrence}. Let \(\N_0\)
be its reference chain dimension. Its exact compact reduced homological
profile can be computed using one ordinary reference solve, with safe
time bound
\begin{equation}\label{eq:reference-time}
\poly(b,W)\,\N_0^3+\poly\bigl(W,\log(m+1),\log(\N_0+1)\bigr).
\end{equation}
The exponent-dependent part of the output has polynomial bit size in
\(W+\log(m+1)+\log(\N_0+1)\).
For original one-component closures with at least one run, this gives the
coarser uniform bound
\begin{equation}\label{eq:context-exp}
2^{O(W+1)}\,\poly(\log(m+1)).
\end{equation}
Fixed resource caps can interrupt the reference solve; in that event the
bounded implementation makes no exact-rank claim.
\end{theorem}

\begin{proof}
The reference has \(2W+2\) expanded crossings. There are at most \(W\)
other nonempty runs, so \(t\le W+1\).
Each reference state can be constructed and its local Frobenius maps
evaluated in a polynomial in \(b,W\), times the number of labels processed.
Assembling the reference differential and applying ordinary elimination
is therefore bounded by the first term of \eqref{eq:reference-time}.
This intentionally loose cubic bound suffices regardless of the practical
bitset improvements.

The reference homological support has at most \(2W+3\) possible degrees.
Splitting, shifting, and summing these entries takes polynomially many
integer operations. All reference coefficients and boundary ranks are at
most \(\N_0\), and the final rank is at most
\(\N_0+m\N_0\). Degree shifts use integers of \(O(\log(m+W+2))\)
bits. Standard integer arithmetic proves the second term and the output
claim.

For the uniform knot bound, every generator occurs in the original braid.
At most one generator index is supplied only by the selected run, so
\(b-2\le W\). A support of size \(s\) has at most \(b+s\) resolution
circles: changing one \(I\) into \(E\) changes component count by at most
one. Hence, even without the sharper first-use bound from report 11,
\[
\N_0\le 2^{b-1}\prod_j(1+2m_j^{(0)})
\le 2^{b-1}3^{\,2W+2}=2^{O(W)}.
\]
Here \(1+2u\le3^u\) for integers \(u\ge1\), and \(m_j^{(0)}\)
are the reference lengths. Substituting this into
\eqref{eq:reference-time} proves \eqref{eq:context-exp}.
The reference may be a link; the argument only uses occurrence of the
original generator indices, which shortening preserves.
\end{proof}

For fixed context the finite solve is a constant with respect to \(m\),
and the remaining bit complexity is polynomial in \(\log m\). This is
the encoding improvement over expanding all \(m+1\) local terms.
In particular \(W=O((\log n)^2)\) gives \(n^{O(\log n)}\), and
\(W=O((\log\lambda)^2)\) gives \(\lambda^{O(\log\lambda)}\).
More generally, if \(W=O((\log n)^k)\) for a fixed \(k\), then
\eqref{eq:context-exp} is a quasi-polynomial homology bound in the
expanded size \(n\). If the same relation is imposed with the encoded
size \(\lambda\), it is a quasi-polynomial bound in \(\lambda\).
These are promises on the supplied context, not theorems that arbitrary
diagrams admit such a representation.

Several comparable huge runs do not satisfy the fixed-small-context
condition. For example, if two lengths are both of order \(m\), choosing
one leaves \(W\) of order \(m\). Applying the theorem once does not remove
that dependence. Any run satisfying \(m\ge W+2\) has more than half the
total crossings, so there can be at most one. Choosing the largest run
therefore finds this regime whenever it is present. This fact is not a
multivariable compression theorem.

\subsection{What can and cannot be inferred from a shortened reference}

The original strand permutation is indispensable. A two-strand run of
length \(101\) is a knot, while its reference length \(m_0=2\) is a
two-component link. Nevertheless the marked reference homology is the
correct algebraic input to the recurrence. Using the reference component
count in a recognition wrapper would be a bug.

The original crossing counts are also indispensable for a grading adapter.
The existing macro-to-scanner conversion uses
\[
h_{\rm cube}=n_+^{\rm original}-h_{\rm macro}.
\]
Substituting the shortened positive crossing count would shift the reported
homology incorrectly. The delivered raw interface retains the macro
grading explicitly. The existing one-component characteristic-two adapter
can double reduced dimensions to obtain its unreduced output, but this
knot-specific convention should not be silently transferred to general
marked links.

The flag \code{check_d2=True} checks the actual finite reference
differential. It does not construct and multiply an enormous original
matrix. The output says which finite complex was checked. General
correctness for all larger exponents rests on the proof of the recurrence.

\section{Structural certificates directly from braid runs}\label{sec:structural}

An exact algebraic improvement is valuable even when recognition can be
settled more cheaply. The audit makes this separation particularly sharp
for the dominant-run family. We first specialize the existing structural
bounds to braid counts, then prove why a proposed signature shortcut
would add no cases.

\subsection{The signed path multigraph}

Assume the original closure is one component. Let \(p_i,n_i\) be the numbers
of positive and negative occurrences of generator \(i\). Define
\[
P=\sum_i p_i,\quad N_-=\sum_i n_i,\quad n=P+N_-,\quad e=P-N_-,
\]
\[
u_+=|\{i:p_i>0\}|,\qquad u_-=|\{i:n_i>0\}|,\qquad d=n-b+1.
\]
The \(b\) oriented Seifert disks of a braid form a path multigraph, with
all crossings of \(\sigma_i\) joining vertices \(i,i+1\). The positive
and negative subgraphs have \(b-u_+\) and \(b-u_-\) connected components,
respectively, including isolated vertices.
Specializing the Kawamura--Lobb bound and the homogeneous-diagram
criterion~\cite{lobb,abe,cromwell} gives
\begin{equation}\label{eq:braid-rasmussen}
e+b-2u_+-1\ \le s(K)\le\ e-b+2u_-+1,
\end{equation}
in the usual Artin crossing-sign convention.
The canonical surface has genus \(d/2\), and the homogeneity defect is
\begin{equation}\label{eq:mixed}
u_++u_--b+1
=|\{i:p_i>0\text{ and }n_i>0\}|.
\end{equation}
The equality follows because every generator index occurs in at least
one sign.

Thus \(d=0\) certifies the unknot. If the defect is zero, the homogeneous
diagram realizes the knot genus; positive \(d\) certifies knottedness.
An interval in \eqref{eq:braid-rasmussen} excluding zero also certifies
knottedness. An interval containing zero alone does not certify triviality,
even if it consists of the single value zero.

\begin{proposition}[Encoded structural evaluation]\label{prop:rle}
All these quantities and the original one-component condition can be
computed in polynomial time in the run-encoded input size \(\lambda\),
without expanding exponents.
\end{proposition}

\begin{proof}
Accumulate signed occurrence counts and exponent sums using integer
arithmetic. Apply the adjacent transposition for each odd exponent, keeping
only touched positions in a sparse permutation. A position not touched is
fixed; add those fixed cycles to the cycles in the sparse permutation.
If some generator index is absent, reject the knot-only interface without
allocating an array of length \(b\).
Otherwise, for a knot candidate \(b-1\le t\), so all count and cycle
structures have polynomial size in \(\lambda\).
The displayed formulas then require only additions, comparisons, and
small-integer multiplications on numbers of \(O(\lambda)\) bits.
\end{proof}

The module \code{braid_profile.py} implements this specialization.
One repository sign convention needs explicit documentation.
For inputs built by \code{Diagram.from_braid}, the production Seifert
evaluator assigns the opposite sign to the usual positive Artin letter.
Consequently the public fields matching existing Seifert output use writhe
\(-e\) and interval \([-U,-L]\), where \([L,U]\) is
\eqref{eq:braid-rasmussen}. The module also returns
\code{artin_exponent_sum} and \code{artin_rasmussen_interval}.
This sign conversion leaves every recognition conclusion unchanged.
Exhaustive comparisons check the complete existing field dictionary,
not just the final verdict.

\subsection{A sound signature idea that is dominated}

The proposed shortcut was to retain many same-sign bands in the canonical
Seifert surface and use a large definite subspace of its symmetrized form.
The mathematical argument is valid, including when selected occurrences
are nonconsecutive in the original word.

\begin{lemma}[Retained bands]\label{lem:bands}
Retain all \(b\) Seifert disks and \(p\ge1\) same-sign occurrences of one
generator. The retained surface contributes a definite subspace of
dimension \(p-1\) to the symmetrized Seifert form of the full surface.
For a fixed sign, nonadjacent generator indices contribute a direct sum
of such subspaces.
\end{lemma}

\begin{proof}
The retained surface is the standard \(T(2,p)\) band surface, together
with unused disks. Its symmetrized form is definite of rank \(p-1\),
with sign determined by the crossing convention~\cite{feller}.
Deleting other bands does not alter the relative order of the retained
same-generator bands.
The full surface is obtained by attaching only one-handles to this
retained surface \(F_0\). Hence \(H_2(F,F_0)=0\), and the relative
homology sequence gives an injection \(H_1(F_0)\to H_1(F)\).
The Seifert pairing restricts to the pairing of \(F_0\).
For nonadjacent indices the retained bands use disjoint strand pairs and
form a split union after deleting all other bands; their forms sum
directly. Retaining all disks is relevant to the injection argument:
arbitrary subsurfaces need not have injective first homology.
\end{proof}

For one sign let \(w_i=\max(p_i-1,0)\), using \(n_i\) for the other
sign. A maximum-weight independent set of the index path chooses
nonadjacent generators with total weight \(R\). It is computed by
\[
B_j=\max(B_{j-1},B_{j-2}+w_j).
\]
A definite \(R\)-dimensional subspace in a symmetric form of dimension
\(d\) gives the sufficient signature inequality
\begin{equation}\label{eq:inertia}
|\sigma(K)|\ge 2R-d
\quad\text{when the right side is positive}.
\end{equation}
Indeed, at least \(R\) eigenvalues have the selected sign and at most
\(d-R\) can have the opposite sign.

\begin{proposition}[Domination]\label{prop:domination}
Every braid knot certified nontrivial by \eqref{eq:inertia} is already
certified by \eqref{eq:braid-rasmussen}. Signature packing adds no
instances to the existing structural stage.
\end{proposition}

\begin{proof}
For positive bands,
\[
R\le\sum_i\max(p_i-1,0)=P-u_+.
\]
The lower endpoint of \eqref{eq:braid-rasmussen} equals
\[
2(P-u_+)-d.
\]
It is at least \(2R-d\). Therefore a positive signature packing bound
forces a positive Rasmussen lower bound. The negative version gives an
upper bound at most \(-[2(N_--u_-)-d]\), and the same conclusion follows
when negative packing succeeds.
\end{proof}

The diagnostic signature routine and witness checker are retained for
reproducibility, but the recognition path uses the cheaper dominating
counts directly. This rejected optimization is a useful result of the
audit: a correct new certificate is not automatically a stronger one.

\begin{corollary}[Every dominant-run knot is already decided]\label{cor:dominant}
If a one-component \(b\)-braid has a selected run of magnitude
\(m\ge W+2\), then its existing structural Rasmussen interval excludes
zero. In particular it is knotted.
\end{corollary}

\begin{proof}
For a positive selected run,
\[
P-u_+\ge m-1,\qquad
2(P-u_+)-d\ge m-W+b-3\ge b-1>0.
\]
Here \(d=m+W-b+1\), and the presence of a run implies \(b\ge2\).
For a negative selected run the upper endpoint is at most the negative
of the same positive quantity.
\end{proof}

This corollary fixes the interpretation of the performance results.
The recurrence computes much more homological information at much lower
cost on its long-run families, but it does not create a new class of
previously undecided knots in that range. The direct run-based structural
interface, on the other hand, makes the already known decision available
without first paying the expanded-input cost.

\section{Degree-streamed exact homology}\label{sec:streaming}

\subsection{Enumeration in the correct signed degree}

The streamed backend evaluates exactly the same macro complex as
\code{twist.core}. It changes storage and evaluation order.
Set \(\nu=\sum_{e_j<0}m_j\) and replace each local coordinate by
\[
\alpha_j=\begin{cases}
k_j,&e_j>0,\\
m_j-k_j,&e_j<0.
\end{cases}
\]
Then every differential increments one \(\alpha_j\), and the total degree
is \(h=\sum_j\alpha_j-\nu\).
For negative runs the support condition is \(\alpha_j\ne m_j\), not
\(\alpha_j\ne0\). This is a necessary detail for mixed-sign inputs.

To enumerate degree \(h\), solve the bounded composition
\(\sum_j\alpha_j=h+\nu\), with \(0\le\alpha_j\le m_j\).
If the remaining sum at position \(j\) is \(R\), and
\(M_{j+1}=\sum_{i>j}m_i\), visit only
\begin{equation}\label{eq:composition}
\max(0,R-M_{j+1})\le\alpha_j\le\min(m_j,R).
\end{equation}
Every visited prefix extends to a solution, since the remaining interval
coordinates realize every integer from zero through their total
capacity. Conversely every valid state obeys these bounds.
The iterative depth-first implementation therefore emits each state
once, uses \(O(t)\) generator memory, and avoids rescanning the entire
Cartesian product in each degree.

If \(s_h\) is the number of states in degree \(h\), there are at most
\(s_h\) feasible prefixes at each depth. The total enumeration cost is
\(O(tS)\), because
\(\sum_hs_h=S\) and \(n+1\le S=\prod_j(m_j+1)\) for nonempty run lists.
The empty list is handled separately.

\subsection{Two adjacent indexes and generated columns}

At degree \(h\), keep the source state index and construct the target
index for degree \(h+1\). An index entry contains the coordinate tuple,
support mask, basis offset, and basis dimension. These two indexes
determine every local edge and its destination offset.
Generate a differential column, immediately reduce it against a
dictionary of pivot columns, and discard it if dependent.
Only pivots are retained for rank computation.

When \(r_h\) is complete, \eqref{eq:betti} finalizes degree \(h\).
Release its source index and pivot dictionary, retain \(r_h\) as a
scalar, and advance. By induction each local map occurs exactly once
in the correct matrix. The degree-wise basis order may differ from the
ordinary backend for negative runs, but this only permutes matrix bases.
Exact elimination therefore gives identical chain dimensions, boundary
ranks, and homology.

Geometry is obtained from the pinned union-find construction.
The geometry cache has an entry cap and a cap on total cached vertex
occurrences. The map cache has an entry cap and a cap on descriptor
integer slots. Both are least-recently-used caches; an evicted support
is not retained in an unbounded history set.

Local maps use compact circle-index descriptors. A merge relabels circles
and sends two \(X\)'s on the same output circle to zero. A split evaluates
\(\Delta(1)\) or \(\Delta(X)\). A dot map stores two circle indices,
cancelling coincident dots and killing multiplication on the marked
\(X\)-circle. This avoids a cached table with one entry for each of
\(2^{c-1}\) labelings. The current source state's outgoing descriptors
also have an explicit active-slot cap.

\subsection{A storage bound that includes auxiliary data}

Let \(V_g=b(t+1)\) be the geometry grid size, let
\(L_s=\max_h(s_h+s_{h+1})\), and let \(C_g,C_m,A_m\) be the actual
peak cached geometry vertex occurrences, cached map slots, and active
map slots. On a word model large enough for coordinates and offsets,
the ordinary streamed computation uses
\begin{equation}\label{eq:stream-space}
\begin{split}
O\biggl(& (t+1)L_s+(n+1)+C_g+C_m+A_m+V_g+t\\
&+\max_h(r_h+4)
\left(1+\left\lceil d_{h+1}/w_{\rm bit}\right\rceil\right)\biggr)
\end{split}
\end{equation}
words, where \(w_{\rm bit}\) is the bitset limb size.
The four-vector allowance covers a yielded column, its evolving reduced
value, a shifted unit bit, and an allocated XOR result.
The \(O(n+1)\) term is the retained full output and scalar degree
statistics. Omitting it would wrongly suggest constant total memory for
a long two-strand explicit-profile computation.

For binary coordinates, replace scalar word counts by their bit lengths.
For example the two indexes have the safe bound
\[
O\!\left(L_s[t\log(n+2)+t+\log(\N+1)]\right)
\]
bits, apart from container overhead. Geometry and descriptor indices use
\(O(\log(V_g+1))\) bits, and support keys use \(O(t)\) bits.
Finite transient geometries and descriptors are included as well as
cached entries.

Without checking mode, the conservative algebraic bit envelope is
\[
\max_h(r_h+4)d_{h+1}.
\]
With \code{check_d2=True}, retain two adjacent matrices and compose them
column by column; the old source state index is unnecessary.
Thus there are still only two state indexes. A safe algebraic envelope
in checking mode is
\begin{equation}\label{eq:d2space}
\max_h\left\{
d_{h-1}d_h+d_hd_{h+1}
+\max\bigl((r_h+4)d_{h+1},\,4d_h+3d_{h+1}\bigr)
\right\}.
\end{equation}
Pivots are released before the composition phase. Matrix and pivot
reference counts are capped separately, so a large list of zero columns
does not evade accounting through its zero bit payload.

The ordinary backend retains all state indexes and all matrices, with
bit envelope \(Q=\sum_hd_hd_{h+1}\). Replacing that sum by adjacent
degree costs can save substantial storage, but the bounds are
mathematical payload bounds, not measured Python heap or process-RSS
ratios. The experiments measure newly traced Python allocations
separately.

\subsection{Time, budgets, and exact early decisions}

Let \(G\) count actual geometry constructions, \(M_{\rm map}\) map
compilations, and \(E_h\) emitted local differential terms in degree
\(h\). Safe bounds are
\[
G\le(t+2)S,\quad M_{\rm map}\le tS,\quad E_h\le2td_h.
\]
Neighbor-tuple construction can cost \(O(t^2S)\). The particular
union-find implementation has a safe elementary \(O(V_g^2)\) construction
bound; the analysis does not assume a stronger bound its code has not
established. In addition to topology and label evaluation, column
assembly and elimination cost at most
\[
O\!\left(\sum_h[E_h+d_h(r_h+1)]
\left(1+\left\lceil d_{h+1}/w_{\rm bit}\right\rceil\right)\right)
\]
word operations. Each pivot XOR lowers the leading set bit, so a column
uses at most \(r_h\) pivot XORs. Optional differential-square checking
adds its actual column-composition work.
Bounded caches can repeat geometry and map work. These facts explain
why a memory improvement need not be a time improvement.

The budget distinguishes total macro states from live states, total
enumerated basis from one layer's basis, and live algebraic payload from
cache storage. Product bounds are checked before expanding states.
Huge exponents and huge strand counts are rejected by arithmetic
guards before large allocations or decimal diagnostic formatting.
The deadlines are cooperative, and the units are structural counts,
not hard process-memory reservations.

After degree \(h\) is finalized, define
\[
R_{\le h}=\sum_{i\le h}(d_i-r_{i-1}-r_i).
\]
Every summand is an actual nonnegative homology dimension. For an
original one-component closure, \(R_{\le h}>1\) therefore proves
\code{KNOTTED}. Later degrees cannot cancel homology already computed.
An early result leaves \code{reduced_rank} unset and reports
\code{reduced_rank_lower_bound=R}.
Its completeness field is \code{homology_complete=False}.
If checking mode is enabled, the checked
prefix is reported and the full-complex check flag remains false.
\code{UNKNOT} requires a complete rank-one calculation.

For a positive two-strand run of odd magnitude at least three, four
states suffice for an early
nontriviality result; three suffice for the negative run. These are
transparent raw-backend demonstrations. Such examples are already
easy for the structural front end and do not demonstrate acceleration
of a filter-resistant recognition workload.
